% BHU PYQ -- Computer Science: Fundamentals of Computing (2024-25 NEP)
\documentclass[11pt,a4paper]{article}
\usepackage[a4paper,top=2.5cm,bottom=2.5cm,left=2.8cm,right=2.8cm,headheight=14pt]{geometry}
\usepackage{amsmath,amssymb}
\usepackage[T1]{fontenc}
\usepackage[utf8]{inputenc}
\usepackage{lmodern,microtype,fancyhdr,enumitem,hyperref,lastpage,parskip}

\hypersetup{colorlinks=true,urlcolor=black,linkcolor=black,pdftitle={CSCMD11: Fundamentals of Computing -- BHU NEP 2024-25}}

\pagestyle{fancy}\fancyhf{}
\renewcommand{\headrulewidth}{0.4pt}\renewcommand{\footrulewidth}{0.4pt}
\fancyhead[L]{\small \textit{Banaras Hindu University}}
\fancyhead[R]{\small \textit{CSCMD11: Fundamentals of Computing (2024-25)}}
\fancyfoot[C]{\small Page \thepage\ of \pageref{LastPage}}

\newlist{parts}{enumerate}{2}
\setlist[parts,1]{label=(\alph*),leftmargin=*,itemsep=4pt,topsep=3pt}
\setlist[parts,2]{label=(\roman*),leftmargin=*,itemsep=2pt,topsep=2pt}

\newcommand{\pts}[1]{\hfill \textit{[#1]}}

\begin{document}

\begin{center}
  {\large\bfseries BANARAS HINDU UNIVERSITY}\\[2pt]
  {\normalsize Under Graduate Multidisciplinary Course Examination, 2024-25 (NEP)}\\[6pt]
  \rule{0.8\linewidth}{0.5pt}\\[6pt]
  {\large\bfseries COMPUTER SCIENCE}\\[4pt]
  {\large\bfseries Paper Code: CSCMD11 : Fundamentals of Computing}\\[4pt]
  {\normalsize Time: 2:30 Hours \hfill Full Marks: 70}
\end{center}
\rule{\linewidth}{0.4pt}
\smallskip

\noindent \textit{Instructions: Attempt any \textbf{five} questions. All questions carry equal marks (14 marks each). Write your Roll No.\ at the top immediately on receipt of this question paper.}

\bigskip

\begin{enumerate}[label=\textbf{\arabic*.},leftmargin=*,itemsep=12pt]

  \item 
  \begin{parts}
    \item Differentiate between computer hardware and software. What are the different types of computer software? Discuss in detail. Also give examples of each of them.\pts{7}
    \item Define what is a computer. What are the different classifications of digital computer? Discuss.\pts{7}
  \end{parts}

  \item 
  \begin{parts}
    \item What do you understand by an operating system? What is its role in the functioning of a computer? Also discuss the different types of Operating system in detail.\pts{7}
    \item Differentiate between compiled and interpreted languages. Also give examples for each of them.\pts{7}
  \end{parts}

  \item 
  \begin{parts}
    \item Discuss the different components of a network communication system. Also discuss the different types of data used for communication over network.\pts{7}
    \item Discuss the different criteria used for judging the quality of a network. Also discuss the different types of networks in detail. Also give examples.\pts{7}
  \end{parts}

  \item 
  \begin{parts}
    \item Discuss different types of network topologies. Also discuss their advantages and disadvantages.\pts{7}
    \item What is Internet? Why is it important? Discuss the basic structure of Internet.\pts{7}
  \end{parts}

  \item 
  \begin{parts}
    \item What do you understand by a data center? What is its role in the Internet? Explain with the help of examples.\pts{7}
    \item What are Internet Service Providers (ISPs)? Also, discuss how last-mile connectivity is provided through the Internet.\pts{7}
  \end{parts}

  \item 
  \begin{parts}
    \item Discuss how addressing is done on the Internet. Also differentiate between IPv4 and IPv6 addressing schemes.\pts{7}
    \item Discuss the structure of World Wide Web in detail.\pts{7}
  \end{parts}

  \item 
  \begin{parts}
    \item Differentiate between algorithm and flowchart. Also discuss the advantages and disadvantages of using a flowchart.\pts{7}
    \item Draw a flowchart to print all even numbers between 1 and 20.\pts{7}
  \end{parts}

\end{enumerate}

\end{document}
